\documentclass[a4paper,12pt]{letter}

\usepackage{amsmath}
\usepackage{amssymb}
\usepackage{nopageno}
\usepackage{amsmath}
\usepackage{enumitem}
% \usepackage{mathtools}
%\usepackage{eqnarray,amsmath}
\usepackage[a4paper, total={7.5in, 11.25in}]{geometry}

\begin{document}

Name: \hrulefill\ Neptun code: \hrulefill

\bigskip

\begin{center}
(KTXFI2EBNF) Physics II. Exam~3 Solutions\\
January~19,~2026
\end{center}

\bigskip
 
\begin{enumerate}[label=\Roman*.]
   
\item Moving Reference Frames. Special Theory of Relativity

\begin{enumerate}[label=\arabic*.]
\setcounter{enumii}{0} % Számláló beállítása

\item  Write down the Lorentz transformation equations used in the special theory of relativity. Specify the exact physical meaning of the symbols used in the formulas, including the reference frame in which they are defined. \hfill (5~points)

\smallskip
\textit{Lecture on October 13, 2025, Slides 18--20.}

\item State the postulates formulated by Albert Einstein for the special theory of relativity. (Pay special attention to the precise and exact formulation of the postulates.) \hfill (5~points)

\smallskip
\textit{Lecture on October 13, 2025, Slide~21.}

\end{enumerate}

% ------------------------------------------------------------------------------------------------------------------------------------

\item Limits of the Classical Conceptual Framework

\begin{enumerate}[label=\arabic*.]
\setcounter{enumii}{2} % Számláló beállítása

\item Formulate the de Broglie principle in your own words. \hfill (5~points)

\smallskip
\textit{Lecture on October 20, 2025, Slide 23.}
  
\item What is the Compton effect, what did it prove, and how was this demonstrated? \hfill (5~points)

\smallskip
\textit{Lecture on October 6, 2025, Slides 29--31.}

\item What is meant by the dual nature of electromagnetic radiation? Describe the meaning of this expression. \hfill (5~points)

\smallskip
\textit{Lecture on October 20, 2025, Slides 23--30.}
  
\end{enumerate}
  
% ------------------------------------------------------------------------------------------------------------------------------------

\item Elements of Quantum Mechanics. Physics of Condensed Matter. Theory of Electrical Conduction

\begin{enumerate}[label=\arabic*.]
\setcounter{enumii}{5} % Számláló beállítása
  
\item What is a forbidden band (band gap)? Define the concept. \hfill (5~points)

\smallskip
\textit{Lecture on October 27, 2025, Slides 18--19, 24--28.}

\item Write down the one-dimensional, time-independent Schrödinger equation using the standard notation. Specify the exact physical meaning of the symbols appearing in the equation. \hfill (5~points)

\smallskip
\textit{Lecture on September 25, 2025 by Szilárd ZSÓKA.}


\end{enumerate}

% ------------------------------------------------------------------------------------------------------------------------------------
  
\item Development of Atomic Theory

\begin{enumerate}[label=\arabic*.]
\setcounter{enumii}{7} % Számláló beállítása

\item State the Pauli exclusion principle. \hfill (5~points)

\smallskip
\textit{Lecture on October 20, 2025, Slide 18.}

\item The meaning of the spin quantum number. What does the spin quantum number represent? Describe in detail what is known about its physical meaning. What values can it take? \hfill (5~points)

\smallskip
\textit{Lecture on October 20, 2025, Slide 17.}

\item Given the element with atomic number 10 in the periodic table, write down its electronic configuration using the standard atomic physics notation. \hfill (5~points)

\smallskip
\textit{Lecture on October 20, 2025, Slides 19--22.}
  
\end{enumerate}

\end{enumerate}

% ------------------------------------------------------------------------------------------------------------------------------------

Grades: 0–25~points: Fail (1), 26–30 points: Pass (2), 31–37 points: Satisfactory (3), 38–45~points: Good~(4), 46–50 points: Excellent (5).

\rightline{Dr.~Gabor FACSKO}
\rightline{\textit{facsko.gabor@uni-obuda.hu}}

\vfill

\end{document}
